% ECONOMICS_PROBLEM: {"id": "C73-24", "title": "Subgame-perfect equilibrium in concurrent reachability games", "jel_code": "C73", "source_id": "OP-0159"}
% FIRST_POSTED: 2026-10-10
% ATTEMPT_STATUS: unresolved
\documentclass[11pt]{article}
\usepackage[a4paper,margin=25mm]{geometry}
\usepackage{amsmath,amssymb,amsthm,mathtools,microtype}
\usepackage[hidelinks]{hyperref}
\newtheorem{theorem}{Theorem}
\newtheorem{lemma}[theorem]{Lemma}
\newtheorem{corollary}[theorem]{Corollary}
\newcommand{\R}{\mathbb R}
\newcommand{\E}{\mathbb E}
\newcommand{\supp}{\operatorname{supp}}
\newcommand{\SPE}{\mathrm{SPE}}
\title{Proof of partial results on exact equilibria\\in concurrent reachability games}
\author{Ji Ho Bae}
\date{10 October 2026}
\hypersetup{pdftitle={Proof of partial results on exact equilibria in concurrent reachability games},pdfauthor={Ji Ho Bae}}
\begin{document}
\maketitle
\begin{abstract}
We prove exact subgame-perfect equilibrium existence in a finite two-player
reachability game when, before either target is reached, one player has
at most one decision state and at most two actions there. The opponent's
actions and the transition graph are unrestricted. The proof uses a
finite-policy perturbation and a return-time argument; deviations may be
arbitrary behavioral functions of the public state history. We also compute
the entire exact equilibrium payoff set of a three-state game. It is
nonclosed, and its points have unbounded, exactly determined public-memory
requirements. These are partial results for C73-24. Existence for every
finite two-player concurrent reachability game remains unresolved here.
\end{abstract}

\section{Question and scope}
Let a finite stochastic game have two players, finite nonempty action sets,
initial state $s_0$, transition law $P$, and targets $F_1,F_2$.
At a public state history the players randomize independently. Their
behavioral strategies use that history alone. A finite history is feasible
if each of its transitions has positive probability under some joint action.
For a feasible prefix $h$, let $u_i^h$ be the probability that the full
path, including $h$, visits $F_i$. Thus $u_i^h=1$ when $h$ already visits
$F_i$. Exact subgame perfection means
\begin{equation}\label{eq:spe}
 u_i^h(\tau_i,\sigma_{-i}|_h)\leq u_i^h(\sigma|_h)
 \quad\text{for every feasible }h,\ i\in\{1,2\},\text{ and behavioral }\tau_i.
\end{equation}
This includes prefixes having probability zero under $\sigma$; each
continuation has its own probability law with the prefix fixed.

The general existence question is explicitly left open in
\cite[Section 1]{RV}; its Section 2 gives the history-sensitive conventions
above. The principal partial result of this note is the following.

\begin{theorem}\label{thm:reach}
Suppose that, among states outside $F_1\cup F_2$, player 1 has exactly one
action except possibly at one state $s_*$, where it has at most two.
Then an exact SPE exists. Its actions depend only on the current state
and which targets have already been visited. There is no restriction on
player 2's finite action sets or on the transition probabilities.
The same assertion holds with the players interchanged.
\end{theorem}

The restriction on player 1 is an additional hypothesis, not a property
of general C73-24 instances. Section~\ref{sec:nonclosed} gives a separate,
fully classified example explaining why a limit of equilibrium payoffs
need not itself be an equilibrium payoff.

\section{Terminal-reward games and two elementary tools}
Consider a finite set $S$ of unresolved states and finitely many terminal
outcomes. At $s$, actions $a,b$ produce an unresolved transition kernel
$K(s,a,b,t)$ with row sum at most one and, otherwise, a terminal outcome
with reward $r_i\in[0,M]$. Write $g_i(s,a,b)$ for its expected immediate
terminal reward. The payoff is the expected terminal reward, with value
zero if termination never occurs. The bound $M<\infty$ is fixed.
Let $V_i^{x,y}$ denote the vector of actual payoffs under stationary
behavioral policies $x,y$. A stationary profile that is a best response
for each player from every state is an SPE, including after off-path
prefixes. No action observations are used.

\begin{lemma}\label{lem:mdp}
In a finite Markov decision process with bounded nonnegative terminal
rewards and zero reward on nontermination, a pure stationary policy is
optimal simultaneously from all states against all behavioral policies.
A statewise convex combination with fixed global weights of stationary
policies that are optimal from all states is also optimal from all states.
\end{lemma}
\begin{proof}
For $0<\lambda<1$, discount a reward received on transition $n$ by
$\lambda^n$. The Bellman map is a contraction in the maximum norm; selecting
a maximizing action at each state gives a pure stationary optimum against
all policies. There are finitely many pure stationary policies. Along
some sequence $\lambda\uparrow1$, one policy $\pi$ is therefore optimal
at every term. For each behavioral competitor and each initial state,
monotone convergence in its discounted optimality inequality proves that
$\pi$ is undiscounted-optimal.

Let $v$ be the optimal value vector and let $z$ be the stated convex
combination. The Bellman equality $v=g_z+K_zv$ holds because it holds
for every component policy. Any closed unresolved recurrent class $C$
of $z$ is closed under each component of positive weight: a nonnegative
transition out of $C$ cannot cancel in a convex combination. Each such
component earns zero on $C$, so its optimal value $v$ is zero there.
Iterating the equality gives
\[
 v=\sum_{k=0}^{n-1}K_z^k g_z+K_z^n v.
\]
In a finite chain, the probability of avoiding both termination and the
closed unresolved recurrent classes tends to zero. Since $v=0$ on those
classes, $K_z^n v\to0$. The sum converges to the actual terminal reward
of $z$, proving optimality.
\end{proof}

\begin{lemma}[Crossing stationary policies]\label{lem:cross}
Every finite two-player terminal-reward game as above has fully mixed
stationary profiles $(x_n,y_n)$ converging to $(x_0,y_0)$ such that, for
all sufficiently large $n$, $x_0$ is an all-state best response to $y_n$
and $y_0$ is an all-state best response to $x_n$.
If $(x_0,y_0)$ terminates almost surely from every state, it is an exact SPE.
\end{lemma}
\begin{proof}
Let $\Pi_i$ be player $i$'s finite set of pure stationary policies. A vector
$w_i\in\Delta(\Pi_i)$ defines the stationary behavior
\[
 T_i(w_i)(a\mid s)=\sum_{\pi\in\Pi_i:\,\pi(s)=a}w_i(\pi).
\]
The policy is not sampled once and remembered: the displayed distribution
is used anew at every history ending at $s$.
Fix $\mu(s)>0$ with $\sum_s\mu(s)=1$, and set
\[
 L_i(\pi,z)=\sum_s\mu(s)V_i^{\pi,z}(s).
\]
For a fully mixed stationary opponent $z$, this is continuous in $z$.
Indeed, the support graph with $\pi$ fixed does not change while all of
$z$'s action probabilities stay positive. Its closed unresolved classes
have value zero, and the other values are obtained from the inverse of
the transient matrix $I-K$.

Put $N=\max_i|\Pi_i|$, fix $0<\delta<1$, and restrict each policy
probability to be at least $\eta=\delta^N/N$.
Against $w_{-i}$, prescribe the correspondence of all such $w_i$ satisfying
\begin{equation}\label{eq:proper}
 w_i(\pi)\leq\delta w_i(\pi')
 \quad\text{whenever }L_i(\pi,T_{-i}(w_{-i}))
                  <L_i(\pi',T_{-i}(w_{-i})).
\end{equation}
It has nonempty compact convex values. To see nonemptiness, order the
$m=|\Pi_i|$ policies from best to worst and assign probabilities
proportional to $1,\delta,\ldots,\delta^{m-1}$, breaking ties arbitrarily.
The smallest is at least $\delta^{m-1}/m\geq\eta$.
The graph is closed: every strict payoff inequality at a limit persists
near that limit. Kakutani's theorem \cite[p.~458]{Kakutani} applies to the
product of the two restricted simplices and gives a fixed point.
This is a finite-policy version of the proper-pair construction in
\cite[Definition 2.3 and Theorem 2.4]{FTV}; the required continuity and
optimality properties have been established here.

Take $\delta_n\downarrow0$ and a convergent subsequence
$w_i^n\to w_i^0$. If $w_i^0(\pi)>0$, then eventually $w_i^n(\pi)$ is
bounded below by a positive constant. Equation~\eqref{eq:proper} rules
out a strictly better pure policy, since it would force
$w_i^n(\pi)\leq\delta_n$. Thus every policy in $\supp(w_i^0)$ maximizes
$L_i$ against $T_{-i}(w_{-i}^n)$ for all large $n$.
By Lemma~\ref{lem:mdp}, one pure policy attains the optimum in every
state. Positive weights $\mu(s)$ imply that any scalar maximizer also
attains that optimum in every state. The convexity assertion of the
same lemma proves the crossing property for
$x_n=T_1(w_1^n)$ and $y_n=T_2(w_2^n)$.

If the limit profile terminates almost surely, $I-K_{x_0,y_0}$ is
invertible, so compliance values are continuous in a neighborhood of it.
For each fixed pure stationary deviation $\pi$, nonnegative rewards
make its actual value lower semicontinuous in the opponent's policy:
it is the increasing supremum of the finite-horizon values, which are
continuous polynomials. Consequently, coordinatewise,
\[
 V_1^{\pi,y_0}\leq\liminf_n V_1^{\pi,y_n}
 \leq\lim_n V_1^{x_0,y_n}=V_1^{x_0,y_0}.
\]
The argument for player 2 uses the other crossing inequality.
Lemma~\ref{lem:mdp} bounds every behavioral deviation by a pure stationary
optimum. Hence neither player can gain after any feasible history.
\end{proof}

\section{One binary decision state}
\begin{theorem}\label{thm:terminal}
In a finite two-player terminal-reward game with bounded nonnegative
rewards, suppose player 1 has one action at every unresolved state
except possibly $s_*$, where it has at most two. An exact stationary
SPE exists against all behavioral public-state-history deviations.
\end{theorem}
\begin{proof}
First solve the states where a player's reward is cooperatively
unattainable. Let $H_i(s)$ be the maximum expected reward when a single
controller chooses both players' actions, and let $Z_i=\{H_i=0\}$.
Lemma~\ref{lem:mdp} supplies a pure stationary cooperative optimum.
Nonnegativity implies that $Z_i$ is closed under every unresolved
transition of every joint action and that no positive $i$-reward can
be earned there. On $Z_1$, use a joint stationary policy maximizing
player 2's reward. Player 1 is indifferent and player 2 attains its
cooperative upper bound, so this is an all-state SPE. Do the symmetric
construction on $Z_2$. On their intersection both rewards are always
zero; policies there may be made identical without changing any value.
This solves the closed set $Z=Z_1\cup Z_2$.

Replace first entry into $z\in Z$ by termination with the equilibrium
reward just obtained. In every remaining state both cooperative values
are strictly positive. Indeed, a finite path to a positive terminal
reward either remains outside $Z$ until termination or first reaches a
state of $Z$ where the relevant cooperative maximum was preserved.
Against any fully mixed stationary opponent, either player therefore
has a strictly positive optimal reward in every remaining state: it
can prescribe its actions along such a finite positive-probability path.
No bound uniform over fully mixed policies is needed.

If no state remains, we are done. If $s_*$ has been removed or player 1
has only one action there, a pure stationary optimum for player 2 solves
the remaining MDP. Otherwise apply Lemma~\ref{lem:cross}. For large $n$,
the crossed pair $(x_n,y_0)$ terminates almost surely from every state.
Player 2 is an exact best response to the fully mixed $x_n$, so its
actual value is strictly positive everywhere. A closed unresolved
recurrent class, whose actual reward is zero, is impossible.

If $(x_0,y_0)$ is also terminating, Lemma~\ref{lem:cross} finishes the
proof. Otherwise every closed unresolved recurrent class of this limit
must contain $s_*$. A class avoiding $s_*$ would remain closed under
$(x_n,y_0)$, since player 1 has no choices within it. Thus there is just
one such class $C$. Moreover, $x_0$ uses a single action $A$ at $s_*$:
if both actions had positive probability, $C$ would again remain closed
under the crossed pair. Call the other action $B$.

Fix $y_0$. Outside $s_*$ the chain has no unresolved closed class
avoiding $s_*$. Hence an excursion starting with either action at
$s_*$ reaches termination or returns to $s_*$ almost surely.
An $A$-excursion stays in $C$ and returns almost surely without any
terminal reward. For a $B$-excursion, let $p$ be the probability of
termination before return and $b_i$ its expected terminal $i$-reward.
Here $p>0$: otherwise neither action could cause termination from
$s_*$, contradicting termination under $(x_n,y_0)$.

Any stationary policy using $B$ with probability $t>0$ on each visit
terminates almost surely from $s_*$. Its value there is $b_i/p$, since
\[
 v_i=t b_i+(1-tp)v_i.
\]
An arbitrary player-1 deviation cannot exceed $b_1/p$ there. To verify
this without a finite-horizon assumption, enumerate its successive
$B$-excursions and let $q_k$ be the probability that its $k$th such
excursion is started. Conditional on any starting history, that
excursion has termination probability $p$ and expected reward $b_1$.
The terminating excursions are disjoint, so
\[
 \E[\text{terminal reward}]=b_1\sum_{k\geq1}q_k
 =\frac{b_1}{p}\,\Pr(\text{termination})\leq\frac{b_1}{p}.
\]
The identities follow by summing nonnegative quantities and also cover
policies that eventually wait forever. Before the first visit to $s_*$
or termination, a path from another state has no player-1 choice.
Thus every stationary policy with $t>0$ is an all-state best response
to $y_0$. In particular $x_n$ is one. Player 2 was already a best
response to $x_n$, proving that $(x_n,y_0)$ is an exact stationary SPE.

Finally restore the solved policies on $Z$. No deviation can return
from $Z$ to its complement, or improve its prescribed boundary reward.
The remaining game's best-response inequalities therefore extend to
the full game from every state and after every feasible prefix.
\end{proof}

\begin{proof}[Proof of Theorem~\ref{thm:reach}]
Record two public bits specifying which targets have been visited.
If both bits are set, both players' rewards are fixed at one.
If just player $i$'s bit is set, that player is indifferent. Choose a
joint pure stationary policy maximizing the other player's reachability
probability. Lemma~\ref{lem:mdp}, applied with a success terminal reward
of one, proves that this policy is an SPE in every such state.
The policy can be implemented by independent pure actions and needs
no information beyond the current state and the bits.

Before either target is reached, replace first entry into the solved
region by its actual equilibrium payoff vector, of the form $(1,v)$,
$(v,1)$, or $(1,1)$. This is a nonnegative terminal-reward game;
paths never reaching either target give $(0,0)$. It has the decision-state
restriction of Theorem~\ref{thm:terminal}. Its stationary SPE, combined
with the solved continuations, satisfies~\eqref{eq:spe}. The bits
retain earlier target visits at every prefix, including off-path ones.
\end{proof}

\section{A nonclosed exact equilibrium payoff set}\label{sec:nonclosed}
Consider three states $s,w,\ell$. States $w$ and $\ell$ are absorbing,
and both players' target is $\{w\}$. At $s$ both choose Wait or Quit.
Wait/Wait returns to $s$. Exactly one Quit goes to $w$ or $\ell$ with
equal probability; two Quits go to $w$ surely. Infinite waiting loses.

\begin{theorem}\label{thm:payoffs}
The exact SPE payoff set from $s$, allowing arbitrary behavioral
public-state-history strategies, is
\begin{equation}\label{eq:set}
 \left\{(u_m,u_m):m\geq1\right\},\qquad u_m=\frac{m+1}{2m}.
\end{equation}
For a deterministic public finite controller realizing $u_m$, at least
$m$ unresolved modes are necessary, and $m$ suffice. The smallest mean
termination time among all SPEs with payoff $u_m$ is
\begin{equation}\label{eq:time}
 \frac{(m+1)(2m+1)}{6m}.
\end{equation}
\end{theorem}
\begin{proof}
Before termination a public history is specified by its length.
Write $p_n,q_n$ for the Quit probabilities at length $n$ and $u_n$
for the common actual continuation payoff in an SPE. Immediate Quit
guarantees at least $1/2$, so $u_n\geq1/2$ at every feasible prefix.
Every unresolved tail has a first stage with a positive Quit probability;
an all-Wait tail would instead have actual value zero.

In fact $u_n>1/2$. If $u_n=1/2$, move past any initial Wait/Wait stages
to the first positive-quitting stage and write its probabilities as
$p,q$ and its next unresolved value as $v\geq1/2$. The Bellman identity gives
\[
 \tfrac12=pq+\tfrac12[p(1-q)+(1-p)q]+(1-p)(1-q)v
 \geq\tfrac12+\tfrac12pq.
\]
Thus exactly one of $p,q$ is positive. The other player could Quit
immediately and gain half that positive probability, a contradiction.

At any positive-quitting stage both probabilities must be positive:
if one were zero, the other's supported Quit action would yield $1/2$,
less than its value. A supported action must have the equilibrium
one-step return, because changing just that action is an admissible
deviation. If either probability is one, both are one: against certain
Quit, Quit gives one and Wait gives $1/2$.

At a stage with $0<p,q<1$, indifference gives
\[
 u=\frac{1+q}{2}=\frac{1+p}{2},\qquad
 \frac{1+p}{2}=\frac p2+(1-p)v.
\]
Hence $p=q$, $v=1/[2(1-p)]$, and $p\leq1/2$ because $v\leq1$.
Wait/Wait stages preserve the value. If $p'$ is the Quit probability
at the next positive-quitting stage, then
\begin{equation}\label{eq:recip}
 p'=\frac p{1-p},\qquad \frac1{p'}=\frac1p-1.
\end{equation}
This includes $p=1/2$, followed by $p'=1$.
At any positive-quitting stage the reciprocal is either one or at
least two. Repeated subtraction of one can stay in this set only if
its initial value is an integer $m$. Otherwise it eventually lies
strictly between one and two. Thus the positive Quit probabilities
are forced to be $1/m,1/(m-1),\ldots,1$, separated only by finite
Wait/Wait intervals. The initial payoff is necessarily $u_m$.

For attainment, use modes $k=m,m-1,\ldots,1$. In mode $k$ each player
independently Quits with probability $1/k$. Every unresolved transition
updates $k$ to $\max(k-1,1)$; this also defines behavior at off-path
prefixes following mode one. For $k\geq2$, the actual Bellman return is
\[
 \frac1k+\left(\frac{k-1}{k}\right)^2u_{k-1}=u_k.
\]
Quit gives $(1+1/k)/2=u_k$, and Wait gives
$1/(2k)+(k-1)u_{k-1}/k=u_k$. In mode one, Quit gives one and Wait
gives $1/2$. Under any unilateral deviation, the other player still
Quits surely by mode one. Termination therefore occurs within $m$
stages, and backward induction proves optimality against every
behavioral deviation after every feasible prefix.

A deterministic public finite controller assigns an action distribution
to its current mode and physical state. Equation~\eqref{eq:recip} forces
$m$ distinct positive Quit probabilities along the surviving histories
of any SPE with payoff $u_m$; these require at least $m$ modes.
The construction attains the bound. In that construction survival
after $k<m$ stages has probability $((m-k)/m)^2$. Summing these survival
probabilities gives
\[
 \E\tau=\sum_{k=0}^{m-1}\left(\frac{m-k}{m}\right)^2
       =\frac{(m+1)(2m+1)}{6m}.
\]
Every other SPE with the same payoff can only insert Wait/Wait stages
between the forced quitting stages. Such delays cannot reduce the
mean termination time. This proves~\eqref{eq:time}.
\end{proof}

In particular $(1/2,1/2)$ is the limit of exact SPE payoffs but is not
itself attained. The constructed policies converge pointwise, at each
fixed unresolved history, to perpetual Wait/Wait, whose actual payoff
is $(0,0)$. This game nevertheless has the stationary SPE Quit/Quit.
The memory lower bound concerns prescribed payoffs, not existence of
some finite-memory equilibrium.

\section{Verification and remaining question}
The existence proof uses actual terminal rewards throughout. Bellman
equalities are supplemented by the zero values on unresolved recurrent
classes in Lemma~\ref{lem:mdp}; an unearned limiting value is never assigned
to infinite play. The crossing construction uses stationary behavioral
randomization, not a privately remembered sample of a stationary policy.
The return-time argument proves optimality against every behavioral
deviation. The reachability reduction retains target visits in the prefix.

Theorem~\ref{thm:terminal} does not cover a player with two decision states
or three actions at its sole decision state. In these cases the
single-delay reduction used in its proof is unavailable. Neither theorem
settles the general existence or nonexistence question of C73-24.
The final status of this contribution is \textbf{PARTIAL; general problem
UNRESOLVED}. No claim of literature priority or independent mathematical
verification is made.

\paragraph{Process.}
This manuscript was generated and checked within a Codex mathematical
research session at the author's request. The workflow comprised direct
proof development, counterexample checks, and comparison with the cited
primary sources. The author requested publication as a partial result.
No Lean formalization or independent referee verification is reported.

\paragraph{Publication.}
The standalone manuscript and PDF are available in the separate
\href{https://github.com/jbaelaw/concurrent-reachability-partial-results/tree/7c08ed6e2eb900855d70d68938de25e3dc303bea}{source repository}, with a
\href{https://github.com/jbaelaw/concurrent-reachability-partial-results/releases/tag/v1.0.0}{versioned release}.
This catalogue contribution is explicitly partial.

\begin{thebibliography}{9}
\bibitem{RV}
S. Rajasekaran and M. Y. Vardi,
\emph{Verifying Equilibria in Finite-Horizon Probabilistic Concurrent Game
Systems}, arXiv:2503.14690v7, 22 April 2026.
Sections 1--2.
\url{https://arxiv.org/html/2503.14690v7}.

\bibitem{FTV}
J. Flesch, F. Thuijsman, and O. J. Vrieze,
Recursive repeated games with absorbing states,
\emph{Mathematics of Operations Research} \textbf{21}(4) (1996), 1016--1022.
Definition 2.3 and Theorem 2.4 concern proper pairs; Theorem 3.1 concerns
stationary approximate equilibria.
\href{https://doi.org/10.1287/moor.21.4.1016}{doi:10.1287/moor.21.4.1016}.
\url{https://dke.maastrichtuniversity.nl/f.thuijsman/recursive%20repeated.pdf}.

\bibitem{Kakutani}
S. Kakutani, A generalization of Brouwer's fixed point theorem,
\emph{Duke Mathematical Journal} \textbf{8}(3) (1941), 457--459.
Compact-convex-set corollary, p.~458.
\href{https://doi.org/10.1215/S0012-7094-41-00838-4}{doi:10.1215/S0012-7094-41-00838-4}.
\end{thebibliography}
\end{document}
